% ARCHIVED at v3.96 (2026-09-24) -- author: 6.4 not academic ("something new"), the corridor projectors a trade-off,
% the general conclusion recalibrated, the HardFlow/DPCC steps baseline in Ch 5. Verbatim, not part of the build.

% ---- Ch5 steps paragraph ----
rates \parencite[Sec.~6.1]{romer2025diffusion}, and the endpoint-projection work reports, on this same
D3IL task, the average number of steps the manipulator takes to reach the target on its constraint-satisfying
episodes as one of its three columns \parencite[Sec.~VII-A]{li2025hardflow}. A plan that arrives in

% ---- \subsection{Vision-Conditioned Alignment}\label{sec:setup:me ----
\subsection{Vision-Conditioned Alignment}\label{sec:setup:metrics:aligning}

% ---- D3IL-avoiding is where this thesis makes its central claim,  ----
D3IL-avoiding is where this thesis makes its central claim, and the claim holds: the average-velocity models

% ---- projection, to overtake the per-step projector of \ac{DPCC}  ----
projection, to overtake the per-step projector of \ac{DPCC} at a lower cost, is met over the hump at twenty
evaluations with a single candidate, $0.2$ against $1.8$ violating steps per flight at $291$ against
$355$\,ms per action; with four candidates

% ---- corridor conclusion RQ2 paragraph ----
Endpoint projection beats the per-step projection of \ac{DPCC} in one place: over the hump at twenty
evaluations with a single candidate, where it leaves fewer violating steps at less time per action, $0.2$
against $1.8$ at $291$ against $355$\,ms, on flights about thirty control steps longer. Under the tilt it
does not: its fewest violating steps, $1.4$ against $2.1$, cost $417$ against $360$\,ms, a trade-off. At
three and five evaluations it is as cheap or cheaper but leaves as many violating steps or more, and at one
and two it has no guiding step.

% ---- and with one another, and whether endpoint projection beats  ----
and with one another, and whether endpoint projection beats the per-step projection of \ac{DPCC}.

% ---- The results fall into the lines the study was built on. On t ----
The results fall into the lines the study was built on. On the two D3IL tasks

% ----   \item[D3IL-avoiding: the central result.] ----
  \item[D3IL-avoiding: the central result.]

% ----   \item[D3IL-aligning: the result again.] ----
  \item[D3IL-aligning: the result again.]

% ----   \item[UAV-pillars: D3IL-avoiding reaffirmed on the quadrot ----
  \item[UAV-pillars: D3IL-avoiding reaffirmed on the quadrotor.]

% ----   \item[UAV-corridor: something new.] The concept of \ac{DPC ----
  \item[UAV-corridor: something new.] The concept of \ac{DPCC} works in three dimensions: the violating steps

% ----   \item[UAV-s-curve: the controller matters.] ----
  \item[UAV-s-curve: the controller matters.]

% ---- Endpoint projection beats the per-step projection of \ac{DPC ----
Endpoint projection beats the per-step projection of \ac{DPCC} where the task is run at a budget that gives
it steps to guide, and has nothing to win where one or two evaluations solve the task.

% ----   \item[D3IL-avoiding, operated at one evaluation.] ----
  \item[D3IL-avoiding, operated at one evaluation.]

% ----   \item[D3IL-aligning, at twenty.] ----
  \item[D3IL-aligning, at twenty.]

% ----   \item[UAV-pillars, at one and two.] ----
  \item[UAV-pillars, at one and two.]

% ----   \item[UAV-s-curve, at one.] ----
  \item[UAV-s-curve, at one.]

% ---- 6.4.2 corridor item ----
  \item[UAV-corridor, from one to twenty.] Endpoint projection wins only over the hump at twenty evaluations
    with one candidate, $0.2$ against $1.8$ violating steps at $291$ against $355$\,ms; under the tilt at
    twenty it leaves fewer violating steps at more time, and at three and five as many or more.

% ---- 6.4.3 reading paragraph (v3.95) ----
Read across the environments, the thesis rests on the two D3IL tasks. On D3IL-avoiding the average-velocity
models Pareto-dominate the diffusion model of \ac{DPCC}, with its own temporal U-Net, in fewer control steps
at a thirtieth of its time per action; on D3IL-aligning analytic average-velocity matching does so again,
closer to the target at a third of the time, and endpoint projection is the better projector on the
constraint there. The quadrotor scenes carry the result to a new plant. UAV-pillars reaffirms D3IL-avoiding
when its plans are flown, and shows that a collision with an obstacle the constraints leave out ends a
quadrotor's flight. UAV-corridor shows that the concept of \ac{DPCC} works in a three-dimensional environment
against constraints no demonstration satisfies, and that on demonstrations as simple as straight lines the
generative model does not matter and the budget does. UAV-s-curve adds the controller: in this framework the
same plans succeed or fail by the controller that flies them, and the cascaded geometric controller used
throughout the thesis is the better of the two. For the second research question the budget decides:
endpoint projection beats the per-step projection of \ac{DPCC} at twenty evaluations, on D3IL-aligning and
over the hump of UAV-corridor, and has nothing to guide where one or two evaluations solve the task.

% ---- % v3.95 (author): the three text tables become key points. ----
% v3.95 (author): the three text tables become key points.
